\documentclass[11pt,article,oneside]{memoir}
\usepackage{amsmath,amsthm}
\usepackage{marginnote}
\usepackage{sidenotes}
\usepackage{epigraph}
\chapterstyle{bianchi}
\setsecnumdepth{subsection}
\settocdepth{section}

\newtheorem{prop}{Proposition}
\newtheorem*{thm*}{Theorem}

\begin{document}
\chapter{Lecture 1: Groups}
\epigraph{Symmetry is what we see at a glance.}{\textit{Blaise Pascal}}

\section{Definitions}
A group is a set with an operation.\marginpar{Margin note here.} Sometimes we write\marginnote{Another note.} it multiplicatively.
\sidenote{A sidenote.}

\begin{prop}\label{prop:id}
The identity is unique.
\end{prop}
\begin{proof}[Sketch]
If $e, e'$ are identities then $e = ee' = e'$.
\end{proof}

\begin{thm*}[Lagrange]
The order of a subgroup divides the order of the group.
\end{thm*}

\begin{proof}
See Proposition~\ref{prop:id}.
\renewcommand{\qedsymbol}{$\square$}
\end{proof}

\section{Examples}
Footnotes with marks: first\footnotemark{} and second\footnotemark{}.
\footnotetext{First note.}
\footnotetext{Second note.}

\begin{center}
\begin{minipage}{0.45\textwidth}
Left side.
\end{minipage}\hfill
\begin{minipage}{0.45\textwidth}
Right side.
\end{minipage}
\end{center}

\noindent\begin{tabular}{@{}p{0.5\textwidth}p{0.5\textwidth}@{}}
Cell one & Cell two
\end{tabular}

\vspace{\baselineskip}
\hrule
\end{document}
